\section*{Abstract}
KhetLens is a mobile application that helps farmers in Punjab identify common leaf diseases of wheat, maize and cotton from a photograph of the plant, taken close up or from a short distance, on an ordinary Android phone and without an internet connection. The diagnosis runs on the device using a compact deep learning model trained on field photographs from several regions and images collected locally with our industry partner. The image result is combined with recent local weather to separate diseases that look alike, and the app tells the farmer when it is not sure instead of guessing. Results, treatment guidance and a disease library are available in English, Urdu and Punjabi with audio. This report presents the project vision, literature review, requirements, methodology and design.
\pagebreak

\section*{Executive Summary}
Wheat, maize and cotton are the main field crops of Punjab, and leaf diseases are one of the most common reasons their yields fall short. Most farmers notice the symptoms themselves but have no quick way to confirm what they are seeing or what to do about it. Extension staff are few, and advice from input dealers is not always neutral. Many farming households, however, have access to an Android phone.

KhetLens puts a diagnosis tool on that phone. A farmer photographs the affected plant, close up or from a short distance; the app identifies the crop and the disease on the device, explains the result in the farmer's language with audio, and shows treatment and prevention guidance that has been reviewed with agronomists from our industry partner, Fatima Fertiliser Company Limited. A disease library and a short history of recent diagnoses are available without any network connection.

We have designed KhetLens around the conditions farmers actually face. It works fully offline, using the network only when available to refresh weather data, send feedback and download updates. It uses the recent weather of the farmer's district to help tell apart conditions that look similar in a photograph, such as rust and heat damage. It checks the photograph for blur and poor lighting, recognises when the photograph shows a crop or disease it does not know, and in those cases asks a few simple picture questions or says plainly that it is not sure. Where the app does give a confident answer, it explains why.

Technically, the app runs a compact MobileNetV3-based classifier compressed to 8-bit integers, with one shared backbone and separate output heads for the crop and for the diseases of each crop. The model is trained from public datasets of field photographs from several regions and from photographs collected locally with our partner. As part of the model-building work we will run a structured study to find how many local photographs are needed before further collection stops improving accuracy. The application is built natively for Android, stores its data locally, and talks to a small backend service for one-time verification, feedback collection and updates.

The report is organised as follows. Chapter 1 introduces the document. Chapter 2 presents the project vision: the problem, our goals, the scope and the stakeholders. Chapter 3 reviews the literature and related applications. Chapter 4 specifies the software requirements, interface and data design. Chapter 5 describes our proposed approach and methodology. Chapter 6 gives the high-level and low-level design.
\pagebreak
